% NichidaiMasterExam_R8_3rd_4
\begin{tcolorbox} %%R8_3rd_4
	$\fbox{4}$
\begin{enumerate}
	\item $\frac{d}{dr}\exp{(-r^2)}$を求めなさい。
	\item $x= r\cos{\theta}, y= r\sin{\theta}$とする。ヤコビアン${\partial(x, y)}/{\partial(r, \theta)}$を求めよ。
	\item $\int_0^\infty \int_0^\infty \exp{(-x^2- y^2)}dxdy$を求めよ。
	\item $\int_0^\infty \int_0^\infty\int_0^\infty \exp{(-x^2- y^2- z^2)}dxdydz$を求めよ。
\end{enumerate}
\end{tcolorbox}

\begin{souanswer} %R8_3rd_4_ans
\begin{enumerate}
	\item 合成関数の微分法を用いて
\begin{equation*}
	\frac{d}{dx}\exp(-r^2)= \exp(-r^2)\cdot \frac{d}{dr}(-r^2)= -2r\exp{(-r^2)}
\end{equation*}
	\item\label{R8_3rd_4(2)} $x= r\cos{\theta}, y= r\sin{\theta}$と変数変換すると、
\begin{equation*}
	\det{\begin{pmatrix}dx/dr& dx/d\theta\\ dy/dr& dyd/\theta\end{pmatrix}}= r
\end{equation*}	
	\item $x= r\cos{\theta}, y= r\sin{\theta}$と変数変換すると、
\begin{itemize}
	\item $\exp{(-x^2- y^2)}= \exp{(-r^2(\cos^2\theta+ \sin^2\theta))}= \exp{-r^2}$
	\item 積分領域は$x\ge 0, y\ge 0$だから$0\le r\le \infty, 0\le \theta\le 2\pi$
	\item ヤコビアンは\zcref{R8_3rd_4(2)}より$r$だから$dxdy= rdrd\theta$
\end{itemize}
	\item \label{R8_3rd_4(4)}
\begin{align*}
	\int_0^\infty \int_0^\infty \exp{(-x^2- y^2)}dxdy
	&= \int_{\theta= 0}^{2\pi}\int_{r= 0}^\infty \exp{(-r^2)}rdrd\theta\\
	&= \int_{r= 0}^\infty \exp{(-r^2)}rdr \int_{\theta= 0}^{2\pi}d\theta\\
	&= \left[-\frac{1}{2}\exp{(-r^2)}\right]_0^\infty \left[\theta\right]_{\theta= 0}^{2\pi}\\
	&= \frac{1}{2}\cdot \frac{\pi}{2}= \frac{\pi}{4}
\end{align*}
	\item 積分は以下のように書き換えることができる。
\begin{align*}
	\int_0^\infty \int_0^\infty\int_0^\infty \exp{(-x^2- y^2- z^2)}dxdydz
	&= \left(\int_0^\infty \exp(-x^2)dx\right)	\left(\int_0^\infty \exp(-y^2)dy\right)\left(\int_0^\infty \exp(-z^2)dz\right)\\
	&= I^3
\end{align*}
	\zcref{R8_3rd_4(4)}より$I> 0$であるから$I= \frac{\sqrt{\pi}}{2}$。よって求める積分の値は$I^3= \left(\frac{\sqrt{\pi}}{2}\right)^3= \frac{\pi\sqrt{\pi}}{8}$
\end{enumerate}
\end{souanswer}